\subsection{Recommendations for dataset releases and evaluations}
\label{sec:recommendations}

\begin{enumerate}
\item \textbf{Release the hierarchy.} Every observation should carry a stable
  observation identifier, its acquisition parent (field, source image, or slide), a
  patient or specimen identifier where ethics and law permit, and, if duplicates were
  audited, a duplicate-component identifier.
\item \textbf{Split at the deepest released unit and publish the split.} Evaluate on
  parent-disjoint (or patient-disjoint) folds and release the exact fold manifest with
  its hash, so that results reported under one dataset name are comparable.
\item \textbf{Report both units.} State observation and parent counts for every split,
  and the fraction of test observations with a training sibling.
\item \textbf{Check the structure before training.} The closed forms of
  Section~\ref{sec:structure} (\texttt{cytofair risk-report}) give the expected straddle,
  sibling rate, diversity ratio, and class--parent dependence of a crop-level split from
  the parent table alone.
\item \textbf{Treat one question as one test.} When the same question is asked of many
  architectures, test it once with a mixed model that treats architectures and splits as
  random effects, and size the study by simulation; report per-architecture contrasts
  descriptively.
\item \textbf{Do not certify a split with a probe, a saliency map, or an augmentation.} Leakage
  inflated accuracy with field identity erased from the features, with identical Grad-CAM maps,
  with backgrounds removed, and with strong colour augmentation. Only a parent-disjoint split
  removes it.
\item \textbf{Keep the evidence.} Save per-sample predictions, final weights, and the
  exact split for every run, so that confusion matrices, calibration, and later
  mechanistic analyses do not require retraining.
\end{enumerate}
